\section{Álgebra Tensorial}

\subsection{Propiedades Universais}

\begin{definition}
    Sejam \(P_1, \ P_2\) propiedades de funções. Dados \(X, \ U\) conjuntos e \(\phi \in \mathcal{F}(X,U)\), diz-se que \((\phi, U)\) é \emph{universal} em \(X\) respeito de \(P_1\) e \(P_2\) se: 
    \[ \forall \psi \in \mathcal{F}(X,A ) \mid  P_1, \ \exists ! \widetilde{\psi} \in \mathcal{F}(U, A )\mid  P_2 \text{ e } \widetilde{\psi} \circ \phi = \psi. \]  
\end{definition}

\ejemplo{ 
    \begin{example}
    \begin{itemize}[left = 3cm]
    %\item O par \((\alpha, V)\) é universal em \(I\) \(\leftrightharpoons\) \(\alpha \) é base de \(V\).
        \item[] \hspace{-2.55cm}{\scriptsize \(\bullet\)} \ \( U = V\) \(\F\)-espaço vetorial. 
        \item \(X = I\) e \(\phi = \alpha : I \to V\).
        \item \(P_1 = \) "\(A=W\) é \(\F\)-espaço vetorial".
        \item \(P_2 = \) "ser linear". 
    \end{itemize}
    \end{example}
    \tcblower
\centering \((\alpha, V)\) é universal em \(I\) \textcolor{verdeoscuro}{\(\leftrightharpoons\)} \(\alpha \) é base de \(V\). \comentario{Olha \cite[Pág. 337]{MA719}}
}

\begin{tikzpicture}[scale = 1.1, overlay, node distance=1.5cm, shift= {(12.5cm, 4.5cm)}]
    % Flecha con límite
    \node (X) {$I$};
    
    % Función psi debajo
    \node (U) [right=of X] {$V$};
    
    % Coma y psi tilde al lado
    \node (A) [below=of X] {$W$};

    \draw[->, color = black] (X) -- node[midway, above] {\footnotesize $\alpha$} (U); 
    \draw[->, color = black] (X) -- node[midway, left] {\footnotesize \textcolor{gray}{$P_1\mid  $ } $\psi$} (A); 

    \draw[->, color = black, dashed] (U) -- node[midway, right, xshift = 0.5cm] {\footnotesize $\widetilde{\psi}$ \textcolor{gray}{$\mid P_2$}} (A); 

    \node at (0.6, -0.6) {\large\rotatebox{180}{$\circlearrowright$}}; 
\end{tikzpicture}

\- \vspace{-2cm}
\hypertarget{lemma1012}{}
\begin{lemma}
    \textbf{10.1.2.} Sejam \((\phi, U)\) universal em \(X\) respeito de \(P_1\) e \(P_2\), \(\psi_1\), \(\psi_2 \in \mathcal{F}(U,A) \mid P_2\). Se \(\psi_1 \circ \phi = \psi_2 \circ \phi \mid P_1\), então \(\psi_1 = \psi_2\).    
\end{lemma}
\demo{
    Segue da universalidade de \((U,\phi)\) em \(X\) tomando \(\psi = \psi_1 \circ \phi\). 
}

\begin{definition}
    Uma propiedade funcional \(P\) é \emph{compatível com composições} se \(f \circ g \mid P\) toda vez que \(\exists f \circ g\) e ambas \(f\text{ e } g\) verifiquem \( P \). 
\end{definition}

\hypertarget{lemma1013}{}
\begin{lemma}
    \textbf{10.1.3.} Sejam \((\phi_a, U)\) e \((\phi_b, V)\) universais em \(X\) respeito de \(P_1\) e \(P_2\). Se \(\phi_a, \  \phi_b \mid P_1\); \ \(\id_{_U}, \ \id_{_V} \mid P_2\) e \(P_2\) compatível com composições, então \(\exists ! f : U\to V\mid P_2\) tal que \(\phi_b = f \circ \phi_a\). Além disso, aquela \(f\) é bijeção. 
\end{lemma}

\demo{
    \parbox{0.21\textwidth}{
        \-
    }
    \parbox{0.78\textwidth}{\vspace{0.2cm}
    Da universalidade \(\exists ! \widetilde{\psi}_a, \ \widetilde{\psi}_b\) como no diagrama. Veja que \(\widetilde{\psi}_a\) e \(\widetilde{\psi}_b\) são inversas a uma da otra. Começe por mostrar primeiro que \((\widetilde{\psi}_a \circ \widetilde{\psi}_b) \circ \phi_a = \id_{_U} \circ \phi_a \) e \((\widetilde{\psi}_b \circ \widetilde{\psi}_a) \circ \phi_b = \id_{_V} \circ \phi_b \), conclua usando o  {\scriptsize \(\square\)} \hyperlink{lemma1012}{10.1.2.}. \vspace{0.3cm}
}
}

\begin{tikzpicture}[overlay, node distance=1.5cm, shift= {(1.15cm, 3.4cm)}]
    % Flecha con límite
    \node (X) {$X$};
    
    % Función psi debajo
    \node (U) [right=of X] {$U$};
    
    % Coma y psi tilde al lado
    \node (A) [below=of X] {$V$};

    \draw[->, color = black] (X) -- node[midway, above] {\footnotesize $\phi_a$} (U); 
    \draw[->, color = black] (X) -- node[midway, left] {\footnotesize $\phi_b$} (A); 

    \draw[->, color = black, dashed] ($(U) + (-0.5cm, -0.3cm)$) -- node[midway, left, xshift = -0cm, yshift = 0.35cm] {\footnotesize $\widetilde{\psi}_a$} ($(A) + (0.3cm, 0.5cm)$); 
    \draw[->, color = black, dashed] ($(A) + (0.5cm, 0.3cm)$) -- node[midway, right, xshift = -0.1cm, yshift = -0.3cm] {\footnotesize $\widetilde{\psi}_b$} ($(U) + (-0.3cm, -0.5cm)$); 

%    \node at (0.7, -0.7) {\large\rotatebox{180}{$\circlearrowright$}}; 
\end{tikzpicture}

\- \vspace{-2cm}
\begin{note}
    Pares universais são únicos a menos de um único isomorfismo. 
\end{note}